% Source: analysis/unbounded_scaling/unbounded_scaling_estimates.{csv,json}
% T = ceil(30 * duration / 76); final pool = 76 * (T - 1) - 3.
% FOV evaluations/chunk = 76 * pool; cumulative counts exclude the first chunk.
\begin{table*}[t]
\centering
\small
\caption{\textbf{A fixed-size read does not bound lookup work.}
Implementation-derived counts for MemCam's exhaustive FOV retrieval at 30 fps
with 76 new frames per chunk. The reader fills 76 context slots, one per target
frame; these need not be distinct historical views. Candidate counts and
per-chunk work are shown for the final chunk. Durations are rounded up to whole
chunks. All rows are calculations, not measured latency; the minute-scale
rows extrapolate the same loop. A B32 archive bounds each query's candidates
by 32, independently of duration.}
\label{tab:lookup-scaling}
\setlength{\tabcolsep}{7pt}
\begin{tabular}{lrrrr}
\toprule
& Context slots & Candidates & \multicolumn{2}{c}{FOV evaluations} \\
\cmidrule(lr){4-5}
Video duration & per chunk & per query & Final chunk & Full rollout \\
\midrule
10 s   & 76 & 225     & 17,100    & 33,972 \\
40 s   & 76 & 1,137   & 86,412    & 689,700 \\
60 s   & 76 & 1,745   & 132,620   & 1,588,932 \\
10 min & 76 & 17,933  & 1,362,908 & 161,477,808 \\
60 min & 76 & 107,993 & 8,207,468 & 5,835,347,868 \\
\midrule
B32, any duration & 76 & $\leq32$ & $\leq2,432$ & $\leq2,432(T-1)$ \\
\bottomrule
\end{tabular}

\vspace{2pt}
\parbox{0.96\textwidth}{\footnotesize
$T$ is the number of generated chunks. Counts cover candidate scoring only,
not denoising, feature extraction, memory updates, or end-to-end runtime.
They describe exhaustive retrieval, not indexed search.}
\end{table*}
